\documentclass[12pt,a4paper]{article}
\usepackage{amsmath,amssymb,mathrsfs,tikz,times,pifont}
\usepackage[T1]{fontenc}
\usepackage{enumitem}
\usepackage{multicol}
\newcommand\circitem[1]{%
\tikz[baseline=(char.base)]{
\node[circle,draw=gray, fill=red!55,
minimum size=1.2em,inner sep=0] (char) {#1};}}
\newcommand\boxitem[1]{%
\tikz[baseline=(char.base)]{
\node[fill=cyan,
minimum size=1.2em,inner sep=0] (char) {#1};}}
\setlist[enumerate,1]{label=\protect\circitem{\arabic*}}
\setlist[enumerate,2]{label=\protect\boxitem{\alph*}}
%%%::::::by chnini ameur :::::::%%%
\everymath{\displaystyle}
\usepackage[left=1cm,right=1cm,top=1cm,bottom=1.7cm]{geometry}
\usepackage[colorlinks=true, linkcolor=blue, urlcolor=blue, citecolor=blue]{hyperref}
\usepackage{array,multirow}
\usepackage[most]{tcolorbox}
\usepackage{varwidth}
\usepackage{float} %pour utiliser l'option [H] qui force l'image à apparaître exactement à l'endroit où elle est placée dans le code.
\tcbuselibrary{skins,hooks}
\usetikzlibrary{patterns}
%%%::::::by chnini ameur :::::::%%%
\newtcolorbox{exa}[2][]{enhanced,breakable,before skip=2mm,after skip=5mm,
colback=yellow!20!white,colframe=black!20!blue,boxrule=0.5mm,
attach boxed title to top left ={xshift=0.6cm,yshift*=1mm-\tcboxedtitleheight},
fonttitle=\bfseries,
title={#2},#1,
% varwidth boxed title*=-3cm,
boxed title style={frame code={
\path[fill=tcbcolback!30!black]
([yshift=-1mm,xshift=-1mm]frame.north west)
arc[start angle=0,end angle=180,radius=1mm]
([yshift=-1mm,xshift=1mm]frame.north east)
arc[start angle=180,end angle=0,radius=1mm];
\path[left color=tcbcolback!60!black,right color = tcbcolback!60!black,
middle color = tcbcolback!80!black]
([xshift=-2mm]frame.north west) -- ([xshift=2mm]frame.north east)
[rounded corners=1mm]-- ([xshift=1mm,yshift=-1mm]frame.north east)
-- (frame.south east) -- (frame.south west)
-- ([xshift=-1mm,yshift=-1mm]frame.north west)
[sharp corners]-- cycle;
},interior engine=empty,
},interior style={top color=yellow!5}}
%%%%%%%%%%%%%%%%%%%%%%%

\usepackage{fancyhdr}
\usepackage{eso-pic}         % Pour ajouter des éléments en arrière-plan
% Commande pour ajouter du texte en arrière-plan
\usepackage{tkz-tab}
\AddToShipoutPicture{
    \AtTextCenter{%
        \makebox[0pt]{\rotatebox{80}{\textcolor[gray]{0.7}{\fontsize{5cm}{5cm}\selectfont PGB}}}
    }
}
\usepackage{lastpage}
\fancyhf{}
\pagestyle{fancy}
\renewcommand{\footrulewidth}{1pt}
\renewcommand{\headrulewidth}{0pt}
\renewcommand{\footruleskip}{10pt}
\fancyfoot[R]{
\color{blue}\ding{45}\ \textbf{2026}
}
\fancyfoot[L]{
\color{blue}\ding{45}\ \textbf{Prof:M. BA}
}
\cfoot{\bf
\thepage /
\pageref{LastPage}}
\begin{document}
\renewcommand{\arraystretch}{1.5}
\renewcommand{\arrayrulewidth}{1.2pt}
\begin{tikzpicture}[overlay,remember picture]
\node[draw=blue,line width=1.2pt,fill=purple,text=blue,inner sep=3mm,rounded corners,pattern=dots]at ([yshift=-2.5cm]current page.north) {\begingroup\setlength{\fboxsep}{0pt}\colorbox{white}{\begin{tabular}{|*1{>{\centering \arraybackslash}p{0.28\textwidth}} |*2{>{\centering \arraybackslash}p{0.2\textwidth}|} *1{>{\centering \arraybackslash}p{0.19\textwidth}|} }
\hline
\multicolumn{3}{|c|}{$\diamond$$\diamond$$\diamond$\ \textbf{Lycée de Dindéfélo}\ $\diamond$$\diamond$$\diamond$ }& \textbf{A.S. : 2025/2026} \\ \hline
\textbf{Matière: Mathématiques}& \textbf{Niveau : T}\textbf{S2} &\textbf{Date: 10/06/2026} & \textbf{Durée : 4 heures} \\ \hline
\multicolumn{4}{|c|}{\parbox[c]{10cm}{\begin{center}
\textbf{{\Large\sffamily Composition Du 2$ ^\text{\bf nd} $ Semestre}}
\end{center}}} \\ \hline
\end{tabular}}\endgroup};
\end{tikzpicture}
\vspace{3cm}
\section*{\underline{Exercice 1 : (5 pts)}}
Soit  $\mathbf{P}$ un polynôme tel que $P(x) = x^3 - 3x + 1$.

On admet que $P(x)$ admet a trois racines distinctes $a$, $b$ et $c$.

\begin{enumerate}
    \item Sans calculer ces racines, déterminer : \textbf{(0,75 $\times$ 6 = 4,5pts)}
\begin{multicols}{3}
\setlength{\columnseprule}{0.1mm} % La largeur de la ligne verticale entre les colonnes
    \begin{enumerate}
    \item $\mathbf{A} = a + b + c$
    \item $\mathbf{B} = ab + bc + ca$
    \item $\mathbf{C} = abc$
    \item $\mathbf{D} = a^2 + b^2 + c^2$
    \item $\mathbf{E} = \dfrac{1}{a} + \dfrac{1}{b} + \dfrac{1}{c}$
    \item $\mathbf{F} = \dfrac{1}{a^2} + \dfrac{1}{b^2} + \dfrac{1}{c^2}$
    \end{enumerate}
\end{multicols} 
    \item En effectuant la division euclidienne de $x^5$ par $x^3 - 3x + 1$, calculer $G = a^5 + b^5 + c^5$. \hfill \textbf{0,5pt}
\end{enumerate}
\section*{\underline{Exercice 2 :(6 pts)} }
\begin{enumerate}
    \item Résoudre dans $\mathbb{R}$ les équations suivants :
    \begin{multicols}{2}
    \setlength{\columnseprule}{0.1mm} % La largeur de la ligne verticale entre les colonnes
    \begin{enumerate}
    \item $3x^2 - 5x + 2 = 0$\hfill \textbf{(1 pt)}
    \item $2|x - 2|^2 + 7|x - 2| + 3 = 0$\hfill \textbf{(1 pt)}
    \item $\left(x + \dfrac{1}{x}\right)^2 = 4\left(x + \dfrac{1}{x}\right) - 3$\hfill \textbf{(1 pt)}
    \end{enumerate}
    \end{multicols}
    \item Résoudre dans $\mathbb{R}$ inéquations suivants :
   \begin{enumerate}
     \item $x^2 - x + 6 \geq 0$\hfill \textbf{(1 pt)}
     \item $\dfrac{2x^2 + 3x - 5}{-x^2 + 3x + 10} \leq 0$\hfill \textbf{(1 pt)}
   \end{enumerate} 
        \item Résoudre dans $\mathbb{R}^2$ les systèmes suivants :
    \begin{multicols}{2}
    \setlength{\columnseprule}{0.1mm} % La largeur de la ligne verticale entre les colonnes 
    \begin{enumerate}
    \item $\begin{cases} x + y = 3 \\ x^2 + 3xy + y^2 = 11 \end{cases}$ \hfill \textbf{(1 pt)} 
    \item $\begin{cases} x^2 + y^2 = 5 \\ xy = 2 \end{cases}$ \hfill \textbf{(1 pt)}
    \item $\begin{cases} x^3 + y^3 = 7 \\ x + y = 1 \end{cases}$ \hfill \textbf{(1 pt)}
    \end{enumerate}
    \end{multicols}
    \item En utilisant la méthode de \textbf{Cramer}, résoudre le système : 
    $\begin{cases} 3x + 4y = 5 \\ 2x - 3y = 1 \end{cases}$ \hfill \textbf{(1 pt)}
\end{enumerate}

\newpage
% ==================== EXERCICE 3 ====================
\section*{\underline{Exercice 3 :(6 pts)} }
Soit $ABC$ un triangle quelconque et $I$ milieu de $[BC]$..

\begin{enumerate}
    \item Déterminer l'ensemble des valeurs réels de $\alpha$ pour que le système suivants admet un barycentre que l'on notera $G_\alpha =\ \{(A, 2\alpha^2 - 3) ; (B, \alpha) ; (C, -\alpha)\}$ \hfill \textbf{(1pt)}
    
    \item Montrer que pour tout réel $\alpha$ on a $\overrightarrow{AG_\alpha} = -\frac{\alpha}{2\alpha^2 - 3}\overrightarrow{BC}$ \hfill \textbf{(1pt)}
    
    \item Construire les points $G_1$ et $G_{-1}$ en utilisant (2) ainsi que le point $I$ \hfill \textbf{(1pt)}
    
    \item Montrer que $I$ est le milieu de $[G_1 ; G_{-1}]$ \hfill \textbf{(1pt)}
    
    \item Déterminer l'ensemble (E) des points M du plan tel que :\hfill \textbf{(1pt)}
    \begin{multicols}{2}
\setlength{\columnseprule}{1mm} % La largeur de la ligne verticale entre les colonnes
    \begin{enumerate}
        \item $\left\|\rule{0pt}{12pt} -\overrightarrow{MA} + \overrightarrow{MB} - \overrightarrow{MC}\right\| = 9$ 
        \item $\left\|\rule{0pt}{12pt} -\overrightarrow{MA} + \overrightarrow{MB} - \overrightarrow{MC}\right\| = \left\|\rule{0pt}{12pt} \overrightarrow{MA} - \overrightarrow{MB}\right\|$ 
    \end{enumerate}
    \end{multicols} 
\end{enumerate}
\end{document}